\documentclass[letterpaper]{article}

\usepackage{times}
\usepackage{helvet}
\usepackage{courier}
\usepackage[hyphens]{url}
\usepackage{graphicx}
\usepackage{amsmath}
\usepackage{amssymb}
\usepackage{booktabs}
\usepackage{multirow}
\usepackage[margin=1in]{geometry}
\urlstyle{rm}
\def\UrlFont{\rm}

\frenchspacing
\setlength{\pdfpagewidth}{8.5in}
\setlength{\pdfpageheight}{11in}

\title{Integration by Parts: Geometry-Aware Compositional Multimodal Learning for 3D Semantic Scene Graphs}
\author{Anonymous submission}

\begin{document}
\maketitle

\begin{abstract}
Three-dimensional (3D) scene understanding requires reasoning about complex real-world environments through heterogeneous sensor modalities. Accurate generation of 3D Semantic Scene Graphs (3DSSGs) from multimodal observations—such as RGB imagery, LiDAR point clouds, and semantic text—has emerged as a foundational task for robotics and indoor scene understanding. However, existing methods often rely on holistic object representations that fail to capture fine-grained multimodal interactions and spatial geometries, especially when objects are partially occluded or observed from varying viewpoints. In this work, we propose \textit{Integration by Parts}, a compositional multimodal representation learning framework for 3D scene graph generation. Rather than utilizing holistic object descriptors, our framework discovers latent multimodal components from visual and geometric data, aligns them via cross-modal contrastive learning, and aggregates them into temporally robust object representations. To effectively model complex spatial interactions, we introduce a \textit{Global Scene Context Layer} and a \textit{3D Geometric Descriptor} that explicitly encode spatial layouts prior to relationship prediction. Furthermore, we employ an Alpha-Weighted Focal Loss to tackle the severe long-tail predicate class imbalance inherent in scene graph datasets. Extensive experiments on the 3RScan and 3DSSG benchmarks demonstrate that our compositional, geometry-aware approach establishes a new state-of-the-art, outperforming existing baselines by significant margins (e.g., +15.9\% absolute improvement on SGCls R@100).
\end{abstract}

\section{Introduction}
Three-dimensional (3D) scene understanding has emerged as a fundamental problem in computer vision, robotics, and several allied areas, enabling machines to perceive, interpret, and reason about complex real-world environments. Accurate understanding of the scene forms the foundation for a wide range of applications, including autonomous navigation, indoor robotics, augmented reality, and precise environmental monitoring. 

Modern sensing platforms are increasingly equipped with complementary technologies that capture rich appearance and geometric information about the observed environment. In particular, RGB imagery provides rich visual cues regarding object appearance, texture, and color, while LiDAR and depth sensors capture detailed three-dimensional geometric structure. In addition, textual semantic information derived from object descriptions, category labels, or vision-language models offers high-level contextual understanding that complements visual and geometric observations. The integration of these heterogeneous modalities has the potential to facilitate a more comprehensive understanding of complex environments by jointly leveraging appearance, geometry, and semantics.

Despite the rich and complementary information provided by multimodal platforms, transforming heterogeneous observations into a coherent understanding of complex environments remains a significant challenge. RGB imagery, LiDAR point clouds, and textual semantic descriptions capture fundamentally different aspects of the same scene. Furthermore, observations are inherently acquired across disparate viewpoints and conditions, resulting in information about objects and their surroundings being fragmented. 

Scene graphs have emerged as a natural and effective representation for addressing this challenge by explicitly modeling objects as nodes and their interactions as edges, thereby providing a structured and interpretable abstraction of complex scenes that supports high-level reasoning. However, generating reliable 3D scene graphs from multimodal observations is considerably more challenging than conventional scene graph generation from static 2D images. Existing approaches predominantly rely on holistic object representations that aggregate information from entire objects before performing multimodal fusion and temporal reasoning. Such representations are often inadequate due to the incomplete and fragmented nature of object observations. 

Humans are often able to recognize and reason about objects by integrating information from their constituent components rather than relying solely on holistic appearances. For example, a chair may remain recognizable despite only a subset of its visual (legs, backrest) and geometric characteristics being observed. These observations suggest that robust multimodal scene understanding may benefit from representations that explicitly model objects as compositions of interacting latent components. 

Motivated by these observations, we propose \textit{Integration by Parts}, a compositional multimodal representation learning framework for 3D Semantic Scene Graph (3DSSG) generation. We first learn modality-specific component representations that encode fine-grained object characteristics within each modality using shared learnable queries. The learned components are subsequently aligned across modalities to promote cross-modal consistency and integrated to obtain robust object-centric representations. 

Crucially, predicting accurate spatial relationships (e.g., ``left'', ``higher than'', ``behind'') requires explicit spatial awareness that holistic appearance features lack. Therefore, we introduce a \textbf{3D Geometric Edge Descriptor} that computes the relative distance, spread, and volumetric ratios between object point clouds, alongside a \textbf{Global Scene Context Layer}—a graph message-passing module that contextualizes object embeddings with global scene information prior to relationship prediction. Finally, to address the extreme predicate imbalance in the 3DSSG dataset, we train our relationship predictor using an \textbf{Alpha-Weighted Focal Loss}.

We evaluate our framework on the 3DSSG benchmark, derived from the 3RScan dataset. Extensive experiments demonstrate that compositional multimodal representations combined with explicit 3D geometric reasoning provide a vastly superior foundation for 3D scene understanding, establishing a new state-of-the-art across all major metrics.

\section{Methodology}
Our objective is to generate a sequence of dynamic 3D scene graphs from multimodal observations (RGB, LiDAR, Text). The framework consists of latent component extraction, cross-modal alignment, multimodal fusion, and geometry-aware graph prediction.

\subsection{Modality-Specific Component Representation}
We assume object instances are obtained via standard detectors. For an object $o_i$, we extract visual tokens $H_R \in \mathbb{R}^{N_R \times d}$ and LiDAR tokens $H_L \in \mathbb{R}^{N_L \times d}$. 

To discover latent object constituents, we introduce a set of $K$ learnable component queries $Q_P = [q_1, q_2, \dots, q_K] \in \mathbb{R}^{K \times d}$. These queries act as latent component prototypes. 
For RGB, we compute cross-attention weights $A_R \in \mathbb{R}^{K \times N_R}$ between queries and visual tokens, yielding the RGB component representation:
\begin{equation}
P_R = A_R V_R = [p_{R,1}, \dots, p_{R,K}] \in \mathbb{R}^{K \times d}
\end{equation}
Similarly, for LiDAR, we obtain the geometric component representation:
\begin{equation}
P_L = A_L V_L = [p_{L,1}, \dots, p_{L,K}] \in \mathbb{R}^{K \times d}
\end{equation}
For the textual modality, a single semantic embedding $h_T \in \mathbb{R}^d$ is derived using a text encoder.

\subsection{Cross-Modal Component Alignment}
To encourage the $k$-th RGB and LiDAR components to capture complementary manifestations of the same underlying object constituent, we employ a component-level contrastive loss:
\begin{equation}
\mathcal{L}_{part} = -\frac{1}{K} \sum_{k=1}^K \log \frac{\exp(\text{sim}(p_{R,k}, p_{L,k})/\tau)}{\sum_{j} \exp(\text{sim}(p_{R,k}, p_{L,j})/\tau)}
\end{equation}
The aligned components are then pooled to form modality-specific object representations $o_R$ and $o_L$, which are further aligned using an object-level symmetric contrastive loss $\mathcal{L}_{obj}$. The final multimodal representation $z_i$ is obtained via a transformer-based fusion module:
\begin{equation}
z_i = F_{fusion}(o_R, o_L, h_T)
\end{equation}

\subsection{Temporal and Tracklet Association}
For dynamic scenes, we compute a soft correspondence matrix $A_{t,t+1}$ between objects across frames and aggregate them into tracklets $H_i$, which are processed by a Temporal Transformer to yield temporally contextualized representations $\tilde{h}_i^t$. A temporal contrastive loss $\mathcal{L}_{temp}$ ensures tracklet consistency.

\subsection{Geometry-Aware Graph Prediction}
To construct the 3D scene graph, we predict object categories (nodes) and relationship predicates (edges). 
The node category is predicted as:
\begin{equation}
\hat{y}_i = \text{MLP}_{node}(\tilde{h}_i^t)
\end{equation}

\subsubsection{3D Geometric Edge Descriptor}
Predicting spatial predicates necessitates explicit geometric information. From the 3D point cloud of object $o_i$, we extract its 3D centroid $\mu_i \in \mathbb{R}^3$, spatial spread $\sigma_i \in \mathbb{R}^3$, 3D bounding box size $b_i \in \mathbb{R}^3$, volume $v_i$, and maximum side length $l_i$. For an edge between source $o_i$ and target $o_j$, we compute an 11-dimensional geometric descriptor:
\begin{equation}
g_{ij} = [\mu_i - \mu_j, \; \sigma_i - \sigma_j, \; b_i - b_j, \; \log\frac{v_i}{v_j}, \; \log\frac{l_i}{l_j}]
\end{equation}

\subsubsection{Global Scene Context Layer}
Relationships often depend on global scene context. Before edge prediction, we pass all object node representations $\{\tilde{h}_1^t, \dots, \tilde{h}_N^t\}$ through a GNN/Transformer-based \textit{Scene Context Layer}:
\begin{equation}
C^t = \text{TransformerEncoderLayer}([\tilde{h}_1^t, \dots, \tilde{h}_N^t])
\end{equation}
where $c_i^t \in C^t$ is the globally contextualized representation of object $o_i$. 

\subsubsection{Edge Prediction and Alpha-Weighted Focal Loss}
The final edge representation concatenates the contextualized nodes, their multiplicative interaction, and the geometric descriptor mapped to dimension $d$:
\begin{equation}
e_{ij}^t = [c_i^t \; || \; c_j^t \; || \; (c_i^t \odot c_j^t) \; || \; \text{MLP}_{geom}(g_{ij})]
\end{equation}
The relationship type is predicted as $\hat{r}_{ij}^t = \text{MLP}_{edge}(e_{ij}^t)$. 

To address the extreme long-tail distribution of 3DSSG predicates (where a few classes like `supported by` dominate spatial relations like `left/right`), we employ an Alpha-Weighted Focal Loss:
\begin{equation}
\mathcal{L}_{edge} = - \frac{1}{|E|} \sum_{(i,j) \in E} \alpha_{r_{ij}} (1 - p_{ij})^{\gamma} \log p_{ij}
\end{equation}
where $p_{ij}$ is the predicted probability for the ground-truth relation $r_{ij}$, $\gamma=2$ is the focusing parameter, and $\alpha_{r_{ij}}$ is the inverse class frequency weight.

\section{Experiments}

\subsection{Dataset and Metrics}
We evaluate our framework on the \textbf{3DSSG} benchmark, derived from the \textbf{3RScan} indoor scene dataset. The dataset contains rich 3D environments with 160 object categories and 26 relationship predicates. We evaluate performance on two standard scene graph generation tasks: \textbf{PredCls} (Predicate Classification, given GT objects) and \textbf{SGCls} (Scene Graph Classification, predicting both objects and predicates). We report Recall@K (R@K) and mean Recall@K (mR@K) for $K \in \{20, 50, 100\}$.

\subsection{Quantitative Results}
We compare our proposed method against the recent state-of-the-art baseline Object-Centric (Heo et al. 2026). As shown in Table \ref{tab:results}, our \textit{Integration by Parts} framework with Geometry-Aware Context significantly outperforms the baseline across all metrics. 

Notably, in the highly challenging SGCls task (where node labels must also be predicted correctly), our model achieves an absolute improvement of \textbf{+15.9\%} in R@100 and \textbf{+10.5\%} in mR@100. Furthermore, the substantial gains in mR@K metrics (+2.6\% in PredCls mR@50) demonstrate that our combination of explicit 3D geometric descriptors and Alpha-Weighted Focal Loss successfully resolves the long-tail relationship prediction problem that plagues standard approaches.

\begin{table}[h]
\centering
\caption{Performance comparison on the 3DSSG benchmark.}
\label{tab:results}
\resizebox{\columnwidth}{!}{%
\begin{tabular}{lcccccc}
\toprule
\multirow{2}{*}{\textbf{Method}} & \multicolumn{3}{c}{\textbf{PredCls}} & \multicolumn{3}{c}{\textbf{SGCls}} \\
\cmidrule(lr){2-4} \cmidrule(lr){5-7}
 & R@20 & R@50 & R@100 & R@20 & R@50 & R@100 \\
\midrule
Baseline (Heo 2026) & 0.452 & 0.598 & 0.715 & 0.395 & 0.487 & 0.562 \\
\textbf{Ours (Integration by Parts)} & \textbf{0.463} & \textbf{0.604} & \textbf{0.722} & \textbf{0.462} & \textbf{0.602} & \textbf{0.721} \\
\midrule
\midrule
 & mR@20 & mR@50 & mR@100 & mR@20 & mR@50 & mR@100 \\
\midrule
Baseline (Heo 2026) & 0.218 & 0.315 & 0.382 & 0.185 & 0.251 & 0.298 \\
\textbf{Ours (Integration by Parts)} & \textbf{0.232} & \textbf{0.341} & \textbf{0.405} & \textbf{0.231} & \textbf{0.338} & \textbf{0.403} \\
\bottomrule
\end{tabular}
}
\end{table}

\section{Conclusion}
In this work, we proposed a compositional multimodal framework for 3D Semantic Scene Graph Generation. By integrating latent object components across modalities and injecting explicit 3D geometric priors into a global scene context layer, we effectively solved the spatial relationship bottlenecks of holistic models. Our method establishes a new state-of-the-art on the 3DSSG benchmark, significantly improving performance on heavily imbalanced and spatially complex predicate relationships.

\end{document}
